{\large\textbf{E1 - Magnetic Pendulum (10 pts)}}

The oscillation frequency of a pendulum can be modified by magnetic forces
between the pendulum and its support. In this experiment you will study
pendulum motion in a combined potential from gravitational and magnetic
interaction terms, using the setup shown in Fig. 3.

\textbf{Equipment (see also Fig. 3)}

A\quad Pendulum body with point-like supports and mirror for angle measurement

B\quad Pendulum tower with hard points to support the pendulum, and laser module
for angle measurement

C\quad Rails to support external magnets

D\quad 2 small dipole magnets to be attached to pendulum body (may be
\textbf{green, red, white} or \textbf{yellow})

E\quad 2 identical external dipole magnets (\textbf{black})

F\quad 2 unknown external magnets F1, F2 (\textbf{blue}, F2 is marked with
\textbf{white} dots at its ends)

G\quad Screen for laser spot for angle measurement

H\quad Stopwatch

I\quad Masking tape, e.g. to fasten pendulum tower to table

J\quad Pencil and ruler

\noindent\fbox{\parbox{\dimexpr\linewidth-2\fboxsep-2\fboxrule\relax}{\bfseries
The magnets are quite strong. Be careful not to hurt yourself or damage the
magnets.

Do not look directly into the laser beam, and turn the laser off when not
needed.

When experimenting with the pendulum, make sure the supporting screws are
resting in the grooves on the pendulum tower.

Feel free to mark the pendulum with your pencil if needed.}}

\textbf{Task E1.1 - Masses (1.0 pts)}

The total mass of the pendulum body with attached small dipole magnets is
$M_\mathrm{pen} + M_\mathrm{mag} = (52.3 \pm 0.2)\,\mathrm{g}$.

Determine both $M_\mathrm{pen}$ and $M_\mathrm{mag}$ as accurately as possible.

\textbf{Task E1.2 - Magnetic dipole moments (4.0 pts)}

With external magnets nearby, the magnetic pendulum moves in a combined
potential formed by gravity and magnetic interaction. The resulting pendulum
frequency $\omega$ can be written as a function of natural frequency
$\omega_1$ and ”magnetic frequency shift” $\omega_\mathrm{mag}$:
\begin{equation}
\omega^2 = \omega_1^2 \pm \omega_\mathrm{mag}^2 \tag{1}
\end{equation}

For the case of two \textbf{black} external dipole magnets, symmetrically
placed at a distance $d$ around the pendulum equilibrium position (see
Fig. 1), and small amplitude oscillations the magnetic frequency shift is:
\begin{equation}
\omega_\mathrm{mag}^2 = \frac{6\mu_0}{I\pi} \cdot j_1 \cdot j_2 \cdot \frac{\ell^2}{d^5}, \tag{2}
\end{equation}
where $\mu_0 = 4\,\pi \cdot 10^{-7}\,\mathrm{N/A^2}$ is the permeability of
vacuum, $I$ is the moment of inertia of the magnetic pendulum around the axis
of rotation, $j_1$ is the combined magnetic moment of the pendulum magnets,
$j_2$ is the magnetic moment of each external dipole, and $\ell$ is the
distance of the pendulum magnet to the rotation axis. For the relative strength
of the dipole moments you may assume $j_2 = 2.4 \cdot j_1$. Local gravity is
$g = 9.81\,\mathrm{m/s^2}$.

\begin{center}\includegraphics[width=0.54\linewidth]{figures/EuPhO_2023_Q4_fig1.png}\end{center}

Figure 1: Frequency shifting using external dipole magnets (top view). $d$
denotes the distance between magnet centers. Note that the orientation of the
external magnets may be reversed.

\textbf{a)} Measure the pendulum frequencies for different magnet distances
$d$, using very small amplitudes. Make sure to cover the whole accessible
frequency range.

\textbf{b)} Determine the ”average magnetization” (magnetic moment per unit
mass) of the material of pendulum magnets and external dipole magnets. Create a
relevant graph for your analysis. Auxiliary measurements may be necessary to
determine all unknowns. You may neglect the mass and thickness of the
non-magnetic coating of the magnets.

\noindent\fbox{\parbox{\dimexpr\linewidth-2\fboxsep-2\fboxrule\relax}{\bfseries
Precise alignment of the rails is important. Make sure that, with the pendulum
in its equilibrium position, the centers of all magnets are on a single line.

Make sure to use symmetric configurations to cancel the force on the pendulum
magnets along the direction of the rails.}}

\textbf{Task E1.3 - Unknown external magnets (3.0 pts)}

The two \textbf{blue} unknown external magnets (F1, F2) each contain several
magnetic dipoles. The dipoles inside F1 are reversed with respect to those
inside F2. The magnetic frequency shift in a setup analogous to Fig. 1 also
follows a power law:
\begin{equation}
\omega_\mathrm{mag,F}^2 \propto d^\alpha. \tag{3}
\end{equation}

\textbf{a)} Measure the pendulum frequencies for different distances $d$,
using very small amplitudes. Choose settings that allow finding the magnetic
frequency shift as accurately as possible.

\textbf{b)} Determine the power law exponent $\alpha$.

\textbf{c)} Sketch a possible configuration of magnetic dipoles inside F1 and
F2 and justify your choice.

\textbf{Task E1.4 - Nonlinear pendulum (2.0 pts)}

Return the setup to the configuration used in Task E1.2, with \textbf{black}
external dipole magnets arranged as in Fig. 1. Following Eqn. 1, the
small-amplitude pendulum frequency can be fully cancelled, $\omega \to 0$.

\textbf{a)} Determine as accurately as possible the magnet separation $d$
required for this full cancellation.

\textbf{b)} Investigate the dependence of pendulum period on its amplitude when
tuned to the best cancellation you were able to obtain.

Suggest a functional dependence and validate it with your data.

Discuss the origin of any possible mismatch.

\textbf{Pictures of experimental setups and equipment}

\begin{center}\includegraphics[width=0.50\linewidth]{figures/EuPhO_2023_Q4_fig2.jpg}\end{center}

Figure 3: Setup and equipment for experimental problem E1.

\noindent\fbox{\parbox{\dimexpr\linewidth-2\fboxsep-2\fboxrule\relax}{\textbf{Note:
The laser module is initially mounted to the setup for E1. You have to take it
off so you can use it for E2. You can also put it back (pay attention to
alignment) when you want to switch back to E1.}}}
